% =============================================
% Section 7: Discussion
% =============================================
\section{Discussion}
\label{sec:discussion}

\subsection{Implications for Edge AI}

The ability to generate coherent natural language on a \$5 microcontroller
has significant implications for edge computing. Unlike cloud-based LLM
inference, on-device generation offers:

\begin{itemize}[leftmargin=*,topsep=2pt,itemsep=1pt]
  \item \textbf{Privacy:} No user data leaves the device. For applications
        such as medical wearables or personal assistants, this eliminates
        data transmission risks entirely.
  \item \textbf{Latency:} Response time is bounded by local computation
        rather than network round-trip time, enabling sub-second interaction.
  \item \textbf{Availability:} The system functions without network
        connectivity, critical for industrial, agricultural, and remote
        deployment scenarios.
  \end{itemize}

\subsection{Hardware Applicability and Portability}

The proposed \tinybit{} framework, through its aggressive quantization to 4-bit (7.5\,MB) or 8-bit (15\,MB), significantly expands the potential for LLM deployment. While our primary target is the ESP32-S3, our analysis indicates suitability for:
\begin{itemize}[leftmargin=*,topsep=2pt,itemsep=1pt]
    \item \textbf{High-end MCUs:} STM32H7/U5 and NXP i.MX RT series equipped with external SDRAM or OctoSPI RAM.
    \item \textbf{RISC-V Ecosystem:} Platforms like the ESP32-P4 (dual-core RISC-V with AI instructions) and Milk-V Duo, leveraging SIMD or custom vector extensions.
    \item \textbf{Memory-Constrained Scenarios:} Devices previously limited to keyword spotting can now host conversational intent engines within <16\,MB footprints.
\end{itemize}

\subsection{Coherence vs.\ Capability Trade-off}

Our coherence analysis reveals a fundamental trade-off: while small models
can achieve grammatical correctness and short-range coherence, long-range
narrative consistency and factual reasoning remain beyond the reach of
sub-15M parameter models. This suggests a natural division of labor:

\begin{itemize}[leftmargin=*,topsep=2pt,itemsep=1pt]
  \item \textbf{On-device:} Template filling, command interpretation,
        short response generation, and local intent understanding.
  \item \textbf{Cloud fallback:} Complex reasoning, multi-turn dialogue,
        and knowledge-intensive tasks.
\end{itemize}

This hybrid architecture, where a small on-device model handles routine
interactions and escalates to a cloud model when complexity exceeds its
capability, represents a practical deployment strategy.

\subsection{Limitations}

We acknowledge several limitations of the current work:

\begin{enumerate}[leftmargin=*,topsep=2pt,itemsep=1pt]
  \item \textbf{Dataset bias:} TinyStories uses simplified vocabulary
        designed for children's stories. Performance on domain-specific
        text (technical, medical, legal) remains untested.
  \item \textbf{Architecture scope:} While we extend our evaluation
        to include a Ternary-Mamba proof-of-concept (Section~\ref{sec:ternary_mamba_eval}),
        more extensive SSM variants and hybrid architectures
        remain unexplored at this scale.
  \item \textbf{Hardware specificity:} Our Xtensa kernel is specific to
        the ESP32-S3. Portability to RISC-V based MCUs (e.g., ESP32-C6)
        requires kernel reimplementation.
  \item \textbf{Evaluation subjectivity:} LLM-as-judge scores, while
        correlated with human judgment, introduce evaluator bias.
\end{enumerate}

\subsection{Quantization-Induced Degradation Analysis}

Our Ternary-Mamba experiments (Section~\ref{sec:ternary_mamba_eval}) reveal an important asymmetry in how extreme quantization affects language generation. While the model's character-level entropy (4.338 vs.\ 4.323 for FP32) remains nearly identical---indicating comparable vocabulary utilization---the 3-gram repetition ratio increases by $\sim$4.7$\times$ (0.019 vs.\ 0.004). This suggests that:

\begin{itemize}[leftmargin=*,topsep=2pt,itemsep=1pt]
  \item \textbf{Syntactic competence is robust to quantization:} The ternary weight space preserves sufficient expressiveness for grammatical structure.
  \item \textbf{Lexical diversity is sensitive to precision:} The coarse weight granularity ($\{-1, 0, +1\}$) reduces the model's ability to disambiguate between competing next-token candidates, leading to repetitive patterns.
  \item \textbf{The Mamba architecture is quantization-friendly:} The SSM's linear-time complexity ($O(N)$ vs.\ $O(N^2)$ for attention) and its recurrent nature may provide inherent advantages for hardware-constrained deployment, as the memory footprint scales linearly with sequence length.
\end{itemize}

\subsection{Future Directions}

Several promising extensions emerge from this work:

\begin{itemize}[leftmargin=*,topsep=2pt,itemsep=1pt]
  \item \textbf{SSM architectures:} Combining MambaLite-Micro's memory
        efficiency with our ternary kernel could push the parameter
        ceiling even higher on the same hardware.
  \item \textbf{Flash-tiered storage:} Using the 16\,MB Flash as a
        secondary model store (analogous to disk-tiered inference)
        could enable models larger than \psram{} capacity through
        layer-by-layer streaming.
  \item \textbf{RISC-V targets:} The ESP32-C6 (RISC-V) and upcoming
        ESP32-P4 (dual-core RISC-V with AI coprocessor) present
        opportunities for ISA-portable kernel development.
  \item \textbf{Task-specific fine-tuning:} A pre-trained TinyBit model
        fine-tuned for specific domains (e.g., Home Assistant voice
        commands) could achieve practical utility even at small scales.
  \item \textbf{On-device learning:} Exploring federated or incremental
        learning to adapt the model post-deployment without cloud
        connectivity.
\end{itemize}
